\exercisesection{分次环的维数}

\begin{enumerate}
\def\labelenumi{\arabic{enumi})}
\tightlist
\item
  设 A 为一个环，p 为 A{[}T{]} 的齐次素理想，且 \(p_{0}\) 为 A 中的素理想 \(p \cap A\)。
\end{enumerate}

\begin{enumerate}
\def\labelenumi{\alph{enumi})}
\item
  若 \(T \in p\)，则 \(p = p_{0} + T \cdot A[T]\)。
\item
  若 T \(\notin\) p，则 \(\mathfrak{p} = \mathfrak{p}_0\)。A{[}T{]}。
\item
  推出在情形 \(a)\) 下有 \(\operatorname{htgr}(\mathfrak{p}) = \operatorname{ht}(\mathfrak{p}_0) + 1\)，而在情形 \(b)\) 下有 \(\operatorname{htgr}(\mathfrak{p}) = \operatorname{ht}(\mathfrak{p}_0)\)。
\item
  证明 \(\dim \operatorname{gr}(\mathbf{A}[\mathbf{T}]) = \dim (\mathbf{A}) + 1\)。
\item
  由此推出一个满足 dim(B) ≠ dimgr(B) 的分次环 B 的例子（p.84, 习题 7）。
\end{enumerate}

\begin{enumerate}
\def\labelenumi{\arabic{enumi})}
\setcounter{enumi}{1}
\tightlist
\item
  设 A 为一个环，\(\mathcal{F} = (\mathcal{F}_i)_{i \in \mathbf{Z}}\) 为 A 的递减滤过，且 \(\mathcal{F}_i = \mathbf{A}\) 对 \(i \leqslant 0\) 成立。令 \(\mathcal{A}\) 为 \(\bigoplus_{i \in \mathbf{Z}} \mathcal{F}_i \cdot v^{-i}\) 在 \(\mathbf{A}[v, v^{-1}]\) 中的子环，令 \(\mathrm{gr}_{\mathcal{F}}(\mathbf{A})\) 为
\end{enumerate}

该滤过环 (A, F) 的伴随分次环。

\begin{enumerate}
\def\labelenumi{\alph{enumi})}
\item
  证明同态 \(f: \mathcal{A} \to \mathrm{gr}_{\mathcal{F}}(\mathbf{A})\) 由 \(f\left(\sum_{i} a_i v^{-i}\right) = \sum_{i} \overline{a}_i\) 定义，其中 \(\overline{a}_i\) 是 \(a_i\) 在 \(\mathcal{F}_i / \mathcal{F}_{i+1}\) 中的类，并诱导出 \(\mathcal{A}/v. \mathcal{A}\) 与 \(\mathrm{gr}_{\mathcal{F}}(\mathbf{A})\) 之间的同构。
\item
  证明对于 A 的每个可逆元素 \(v_{0}\)，由下式定义的同态 \(e_{v_{0}}: \mathcal{A} \to \mathbf{A}\)
\end{enumerate}

\[
e _ {v _ {0}} (\sum a _ {i} v ^ {- i}) = \sum a _ {i} v _ {0} ^ {- i}
\]

诱导出 \(\mathcal{A}/(v - v_{0})\mathcal{A}\) 与 A 之间的同构。

\begin{enumerate}
\def\labelenumi{\alph{enumi})}
\setcounter{enumi}{2}
\item
  证明 \(\mathcal{A}/\mathrm{A}[v]\) 是挠 \(\mathrm{A}[v]\)-模。
\item
  现在假设 A 含有域 k，证明复合态射
\end{enumerate}

\[
k [ v ] \rightarrow \mathrm{A} [ v ] \rightarrow \mathcal {A} \quad (\text {where } \mathrm{A} [ v ] = \bigoplus_ {i \leqslant 0} \mathcal {F} _ {i}. v ^ {- i})
\]

使 \(\mathcal{A}\) 成为无挠的 k{[}v{]}-模。推出 \(\mathcal{A}\) 是平坦的 k{[}v{]}-模。

\begin{enumerate}
\def\labelenumi{\alph{enumi})}
\setcounter{enumi}{4}
\tightlist
\item
  假设 A 是局部环，极大理想为 m，并含有域 k，使得 A（作为加法群）是 k 与 m 的直和。证明 A 存在理想 S，使得 \(\mathcal{A}\) 是 k{[}v{]} 与 S 的直和（我们置 \(\mathcal{F}_{i} = m^{i}\) 对每个 \(i \in \mathbf{Z}\) 成立）。
\end{enumerate}

\begin{enumerate}
\def\labelenumi{\arabic{enumi})}
\setcounter{enumi}{2}
\tightlist
\item
  继续使用上一习题的记号。设 M 为一个 A-模，带有滤过 \((\mathbf{K}_i)_{i\in\mathbf{Z}}\)，使得 \(\mathbf{K}_i = \mathbf{M}\) 对 \(i \leqslant 0\) 成立。
\end{enumerate}

\begin{enumerate}
\def\labelenumi{\alph{enumi})}
\tightlist
\item
  证明分次 A-模 \(\mathcal{M} = \sum_{i\in \mathbf{Z}}\mathrm{K}_{i}v^{-i}\) 具有自然的 \(\mathcal{A}\)-分次模结构。
\end{enumerate}



\begin{enumerate}
\def\labelenumi{\alph{enumi})}
\setcounter{enumi}{1}
\item
  证明分次 A-代数 \(\mathcal{A}\) 由其次数为 1 和 -1 的元素生成，当且仅当 A 存在理想 q 使得 \(\mathcal{F}_i = \mathfrak{q}^i\)（\(i \in \mathbf{Z}\)）。
\item
  假设最后一个条件成立，证明滤过 \((\mathbf{K}_i)_{i\in \mathbb{Z}}\) 是 q-良好的，当且仅当 \(\mathcal{A}\)-分次模 \(\mathcal{M}\) 是有限型的。
\item
  证明若滤过 \((\mathbf{K}_{i})_{i\in\mathbf{Z}}\) 是 q-良好的，则 A-模 M 是 v-无挠的；反之，给定一个无 v-挠的有限生成 A-模 M，可以在其上赋予一个 q-良好滤过，使得伴随 A-模为 M。
\end{enumerate}

\begin{enumerate}
\def\labelenumi{\arabic{enumi})}
\setcounter{enumi}{3}
\tightlist
\item
  采用习题 2 的记号和假设。称滤过 \(\mathcal{F}\) 为准幂滤过，若存在整数 \(k > 0\) 使得
\end{enumerate}

\[
\mathcal {F} _ {n} = \sum_ {i = 1} ^ {k} \mathcal {F} _ {i}. \mathcal {F} _ {n - i} \quad \text { for } \quad n \geqslant 1.
\]

证明若环 A 是诺特环，则下列条件等价：

\begin{enumerate}
\def\labelenumi{(\roman{enumi})}
\item
  滤过 \(\mathcal{F}\) 是准幂滤过。
\item
  环 \(\mathcal{A}\) 是诺特环。
\item
  理想 \(\mathcal{N} = \sum_{i > 1} v^{-i}\mathcal{F}_i\) 是 \(\mathcal{A}\) 的有限型理想。
\item
  存在整数 \(k \geqslant 1\)，使得 \(\mathcal{F}\) 是 A 的关于序关系“\(\mathcal{H}_n \subset \mathfrak{G}_n\) 对每个 \(n\) 都成立”的、以 \(\mathcal{F}_1, ..., \mathcal{F}_k\) 为首项的最小滤过。
\item
  A-代数 \(\mathcal{A}\) 是有限型的。
\end{enumerate}

\begin{enumerate}
\def\labelenumi{\arabic{enumi})}
\setcounter{enumi}{4}
\tightlist
\item
  设 \(\rho: A \to B\) 为诺特环的单同态，使 \(B\) 成为有限型 \(A\)-代数。称 \(B\) 在 \(A\) 上的超越次数为 \(d_A(B)\)，其中 \(d\) 是满足存在 \(A\)-代数单同态 \(b: A[X_1, \ldots, X_d] \to B\) 的最大整数。a) 证明若 \(\sigma: B \to C\) 满足与 \(\rho\) 相同的假设，则有不等式
\end{enumerate}

\[
\mathrm{d} _ {\mathbf {A}} (\mathbf {C}) \geqslant \mathrm{d} _ {\mathbf {A}} (\mathbf {B}) + \mathrm{d} _ {\mathbf {B}} (\mathbf {C}).
\]

\begin{enumerate}
\def\labelenumi{\alph{enumi})}
\setcounter{enumi}{1}
\item
  证明若 A、B 是整环，则 \(d_{\mathbf{A}}(\mathbf{B})\) 等于 B 的分式域在 A 的分式域上的超越次数。此外，若在 a) 中假设 C 是整环，则有等式 \(d_{\mathbf{A}}(\mathbf{C}) = d_{\mathbf{A}}(\mathbf{B}) + d_{\mathbf{B}}(\mathbf{C})\)。
\item
  证明若 B 是 A{[}v, \(v^{-1}\) {]} 的一个包含 A{[}v{]} 的子代数，其中 v 是不定元，则有 \(d_{\mathbf{A}}(\mathbf{B}) = 1\)。
\end{enumerate}

  6) 设 \(\rho: A \to B\) 为诺特环的单同态，使 \(B\) 成为有限型 \(A\)-代数。设 \(q\) 为 \(B\) 中异于 \(B\) 的素理想，并置 \(p = \rho^{-1}(q)\)。a) 证明使用上一习题的记号，有不等式

\[
(*) \quad \mathrm{ht} (\mathfrak {q}) + \mathrm{d} _ {\mathrm{A} / \mathfrak {p}} (\mathrm{B} / \mathfrak {q}) \leqslant \mathrm{ht} (\mathfrak {p}) + \mathrm{d} _ {\mathrm{A}} (\mathrm{B}).
\]

（可以对 \(d_{A}(B)\) 使用归纳法。）

\begin{enumerate}
\def\labelenumi{\alph{enumi})}
\setcounter{enumi}{1}
\item
  证明若 A-代数 B 由 d 个元素生成且 \(\mathrm{d}_{\mathbf{A}}(\mathbf{B}) = d\)，则 (*) 中为等式。
\item
   证明若 A 是整环，且对于每个 n，多项式环 A{[}T \(_{1}\) , \ldots, T \(_{n}\) {]} 都是串联环，则对于每个有限型整 A-代数 B 及 B 中异于 B 的每个素理想 q，(*) 中为等式。
\end{enumerate}

\begin{enumerate}
\def\labelenumi{\arabic{enumi})}
\setcounter{enumi}{6}
\tightlist
\item
  设 A 为诺特环，\(\mathcal{F} = (\mathcal{F}_i)_{i\in \mathbf{Z}}\) 为 A 的递减滤过，使得 \(\mathcal{F}_0 = \mathbf{A}\)、\(\mathcal{F}_1 \neq \mathbf{A}\)，且 A-代数 \(\mathbf{B} = \sum_{i=1}^{\infty} \mathcal{F}_i v^{-i}\) 是有限型的。
\end{enumerate}

我们拟证明 \(\dim(B) = \dim(A) + 1\)，并利用这一事实推广 § 6, no. 3, 命题 5 的推论。将使用前面两个习题的结果。

\begin{enumerate}
\def\labelenumi{\alph{enumi})}
\tightlist
\item
  设 M 为 B 的极大理想。证明有
\end{enumerate}

\[
\operatorname{ht} (\mathfrak {M}) \leqslant \dim (\mathrm{A}) + 1,
\]

并由此推出不等式

\[
\dim (\mathbf {B}) \leqslant \dim (\mathbf {A}) + 1.
\]

\begin{enumerate}
\def\labelenumi{\alph{enumi})}
\setcounter{enumi}{1}
\tightlist
\item
  利用 A{[}v{]}-模 A{[}v, \(v^{-1}\) {]}/A{[}v{]} 的每个元素都被 v 的某个幂消去这一事实，证明有
\end{enumerate}

\[
\dim (\mathrm{A} [ v, v ^ {- 1} ]) = \dim (\mathrm{A} [ v ]) = \dim (\mathrm{A}) + 1,
\]

并推出 \(\dim(\mathbf{B}) \geqslant \dim(\mathbf{A}) + 1\)，从而最终有 \(\dim(\mathbf{B}) = \dim(\mathbf{A}) + 1\)（可以使用环 \(\mathbf{A}[v, v^{-1}]\) 是 \(\mathbf{B}\) 的分式环这一事实）。c) 证明 \(\mathbf{B}\) 的极小素理想是如下形式的理想：

\[
\sum_ {i \in \mathbf {Z}} \left(\mathfrak {p} \cap \mathcal {F} _ {i}\right) v ^ {- i}
\]

其中 p 是 A 的极小素理想。

\begin{enumerate}
\def\labelenumi{\alph{enumi})}
\setcounter{enumi}{3}
\item
  证明 \(v\mathbf{B} \neq \mathbf{B}\)，且 \(v\) 不属于 B 的任何极小素理想。推出不等式 \(\dim(\mathbf{B}/v\mathbf{B}) \leqslant \dim(\mathbf{A})\)（可以使用 § 1 及 § 3, no. 3, 命题 4 的推论 1）。
\item
  现在假设 A 存在包含 \(\mathcal{F}_1\) 的极大理想 m，且 ht(m) = dim(A)。证明在此情形下有 dim(B/vB) ≥ dim(A)，因而 dim(B/vB) = dim(A)。（可以归约到 A 为极大理想 m 的局部环的情形，证明由 v 和 \(\sum_{i\in\mathbf{Z}}(\mathcal{F}_i\cap m)v^{-i}\) 生成的 B 的理想 M 是高度为 dim(A) + 1 的 B 的极大理想，并在 M 处局部化 B。）
\end{enumerate}





\begin{enumerate}
\def\labelenumi{\alph{enumi})}
\setcounter{enumi}{5}
\tightlist
\item
  若 A 为诺特局部环，且 F 为 A 的递减滤过，使得 \(F_{0} = A\)、\(F_{1} \neq A\)，则有等式
\end{enumerate}

\[
\dim (\operatorname{gr} (A)) = \dim (A).
\]

\begin{enumerate}
\def\labelenumi{\arabic{enumi})}
\setcounter{enumi}{7}
\tightlist
\item
  设 A 为诺特环，且 B 为多项式环 A{[}T{]} 的分次子-A-代数。
\end{enumerate}

\begin{enumerate}
\def\labelenumi{\alph{enumi})}
\tightlist
\item
  证明有不等式
\end{enumerate}

\[
\dim (\mathbf {A}) \leqslant \dim (\mathbf {B}) \leqslant \dim (\mathbf {A}) + 1.
\]

\begin{enumerate}
\def\labelenumi{\alph{enumi})}
\setcounter{enumi}{1}
\tightlist
\item
  设 \(B_{i}\) 为 \(b \in A\) 中满足 \(bT^{i} \in B\) 的元素的集合。证明若有 \(\dim(B) = \dim(A)\)，则存在正整数 k，使得对于所有 \(i \geqslant k\)，\(B_{i}\) 都包含于 A 的极小素理想的交中。
\end{enumerate}

\begin{enumerate}
\def\labelenumi{\arabic{enumi})}
\setcounter{enumi}{8}
\item
  设 A 为一个环，M 为有限表示的 A-模。对于每个 \(\mathfrak{p} \in \operatorname{Spec}(A)\)，置 \(\kappa(\mathfrak{p}) = A_{\mathfrak{p}} / \mathfrak{p} \cdot A_{\mathfrak{p}}\)。证明对于所有 \(k \in \mathbf{N}\)，集合 \(\mathfrak{p} \in \operatorname{Spec}(A)\) 中满足 \(\dim_{\kappa(\mathfrak{p})} (M \otimes_A \kappa(\mathfrak{p})) \geqslant k\) 的元素构成 \(\operatorname{Spec}(A)\) 的闭子集。（可以取 M 的表示 \(\mathrm{A}^m \xrightarrow{u} \mathrm{A}^n \longrightarrow \mathrm{M} \longrightarrow 0\)，并考察矩阵 \(u\) 的子式。）
\item
  设 \(\rho: A \to B\) 为环同态。对于每个 \(\mathfrak{p} \in \operatorname{Spec}(B)\)，置
\end{enumerate}

\[
\mathrm{d} (\mathfrak {p}) = \dim_ {\mathfrak {p}} \left(\mathbf {B} \otimes_ {\mathbf {A}} \kappa \left(\rho^ {- 1} (\mathfrak {p})\right)\right).
\]

  本习题旨在证明，当 B 是有限生成的 A-代数时，映射 \(p \mapsto d(p)\) 从 \(\operatorname{Spec}(B)\) 到 \(\mathbf{R}\) 是上半连续的（“Chevalley 定理”）。

记 \(\delta \subset \mathbf{B} \otimes_{\mathbf{A}} \mathbf{B}\) 为典范 A-同态 \(\mathbf{B} \otimes_{\mathbf{A}} \mathbf{B} \to \mathbf{B}\) 的核。对于每个整数 \(r\)，置

\[
\mathrm{P} _ {\mathbf {B} / \mathbf {A}} ^ {r} = (\mathbf {B} \otimes_ {\mathbf {A}} \mathbf {B}) / \delta^ {r}.
\]

称 B ⊗\_A B-代数 P \(_{B/A}^{r}\) 为 r 阶主部代数。通过同态 b ↦ b ⊗ 1 将 B ⊗\_A B 看作 B-代数，于是对于每个 r，P \(_{B/A}^{r}\) 是 B-代数，且典范商同态 P \(_{B/A}^{r+1}\) → P \(_{B/A}^{r}\) 定义出 B-代数的射影系统。

\(^{1}\) 这是作为 \(\kappa(\mathfrak{p})\) 上向量空间的维数。

\begin{enumerate}
\def\labelenumi{\alph{enumi})}
\tightlist
\item
  设 \(p: \mathbf{A} \to \mathbf{A}'\) 为环同态。置 \(\mathbf{B}' = \mathbf{B} \otimes_{\mathbf{A}} \mathbf{A}'\)，并以 \(q: \mathbf{B} \to \mathbf{B}'\) 表示典范同态。设 \(r\) 为整数。证明唯一的 \(\mathbf{B}'\)-代数同态
\end{enumerate}

\[
\mathbf {P} _ {\mathbf {B} / \mathbf {A}} ^ {r} \otimes_ {\mathbf {B}} \mathbf {B} ^ {\prime} \rightarrow \mathbf {P} _ {\mathbf {B} ^ {\prime} / \mathbf {A} ^ {\prime}} ^ {r}
\]

它将 \((b_{1} \otimes b_{2}) \otimes 1\) 的类（\(b_{i} \in \mathbf{B}\)）对应到 \(q(b_{1}) \otimes q(b_{2})\) 的类，是一个同构。b) 设 S 为 B 的一个乘法子集。证明典范同态

\[
\mathrm{S} ^ {- 1} \mathrm{P} _ {\mathrm{B/A}} ^ {r} \rightarrow \mathrm{P} _ {\mathrm{S} ^ {- 1} \mathrm{B/A}} ^ {r}
\]

是一个同构。

\begin{enumerate}
\def\labelenumi{\alph{enumi})}
\setcounter{enumi}{2}
\item
  假设 B 是有限表示的 A-代数，即同构于 \(\mathrm{A}[\mathrm{T}_{1},\ldots,\mathrm{T}_{n}]\) 对某个有限生成理想的商。证明对于每个 \(r\in\mathbf{N}\)，\(\mathbf{P}_{\mathrm{B}/\mathrm{A}}^{r}\) 是有限表示的 B-模。对于每个 \(\mathfrak{p}\in\operatorname{Spec}(\mathbf{B})\)，置 \(\mathrm{P}_{\mathrm{B}/\mathrm{A}}^{\infty}[\mathfrak{p}]=\lim_{r}\mathrm{P}_{\mathrm{B}/\mathrm{A}}^{r}\otimes_{\mathrm{B}}\kappa(\mathfrak{p})\)。证明 \(\mathrm{P}_{\mathrm{B}/\mathrm{A}}^{\infty}[\mathfrak{p}]\) 是一个 \(\kappa(\mathfrak{p})\)-代数，它是局部且完备的，剩余域为 \(\kappa(\mathfrak{p})\)。令 \(k\in\mathbf{N}\)。证明集合 \(\mathfrak{p}\in\operatorname{Spec}(\mathbf{B})\) 中满足 \(\dim(\mathrm{P}_{\mathrm{B}/\mathrm{A}}^{\infty}[\mathfrak{p}])\geqslant k\) 的元素构成 \(\operatorname{Spec}(\mathbf{B})\) 的闭子集。（对于每个 \(r\in\mathbf{N}\)，令 \(\mathrm{F}_{r,k}\) 为 \(\mathfrak{p}\) 中满足 \(\dim_{\kappa(\mathfrak{p})}(\mathrm{P}_{\mathrm{B}/\mathrm{A}}^{r}\otimes\kappa(\mathfrak{p}))\geqslant\binom{r+k-1}{r-1}\) 的元素集合。我们将证明 \(\mathrm{F}_{r,k}\) 是闭集（习题 9），从而 \(\Phi_{k}=\bigcap_{r}\mathrm{F}_{r,k}\) 是闭集。然后应用 § 6, no. 3, 命题 6 的推论 2。）
\item
  设 m 为 B 的一个理想，\(\mu\subset(\mathrm{B/m})\otimes_{\mathrm{A}}\mathrm{B}\) 为典范映射 \((\mathrm{B/m})\otimes_{\mathrm{A}}\mathrm{B}\to\mathrm{B/m}\) 的核。证明有典范同构
\end{enumerate}

\[
\mathrm{P} _ {\mathrm{B/A}} ^ {r} \otimes_ {\mathrm{B}} (\mathrm{B/m}) \rightarrow ((\mathrm{B/m}) \otimes_ {\mathrm{A}} \mathrm{B}) / \mu^ {r}.
\]

假设 A 为域，B 为有限生成的 A-代数，且 m 为 B 的极大理想。证明 \(P_{B/A}^{\infty}[m]\) 同构于 \((B/m)\otimes_{A}B\) 在极大理想 \(\mu\) 处的完备局部化。推出 \(\dim(P_{B/A}^{\infty}[m])=\dim_{m}(B)\)。（应用零点定理和 § 2, no. 3 的定理 1。注意 A 的扩张 B/m 包含于 A 的一个有限纯不可分扩张的有限伽罗瓦扩张中。）

\begin{enumerate}
\def\labelenumi{\alph{enumi})}
\setcounter{enumi}{4}
\tightlist
\item
  假设 B 为有限生成的 A-代数。设 \(p \in \operatorname{Spec}(B)\)。证明有
\end{enumerate}

\[
\dim (\mathrm{P} _ {\mathrm{B/A}} ^ {\infty} [ \mathfrak {p} ]) = \dim \left(\mathrm{B} \otimes_ {\mathrm{A}} \kappa (\rho^ {- 1} (\mathfrak {p}))\right).
\]

  首先注意对于每个 \(r \in \mathbf{N}\)，有 \(\mathbf{P}_{\mathbf{B}/\mathbf{A}}^{\mathbf{r}} \otimes_{\mathbf{B}} \kappa(\mathfrak{p}) = (\mathbf{P}_{\mathbf{B}/\mathbf{A}}^{\mathbf{r}} \otimes_{\mathbf{B}} \mathbf{B}') \otimes_{\mathbf{B}'} \kappa(\mathfrak{p})\)，这是通过置 \(\mathbf{B}' = \mathbf{B} \otimes_{\mathbf{A}} \kappa(\rho^{-1}(\mathfrak{p}))\) 得到的，因而根据 \(a\) )，有 \(\mathbf{P}_{\mathbf{B}/\mathbf{A}}^{\infty}[\mathfrak{p}] = \mathbf{P}_{\mathbf{B}' / \kappa(\rho^{-1}(\mathfrak{p}))}^{\infty}[\mathfrak{p}\mathbf{B}']\)。因此归约到 A 为域时证明等式

\[
\dim \left(P _ {B / A} ^ {\infty} [ p ]\right) = \dim_ {p} (B).
\]

注意当 p 为极大理想时，此等式由 d) 得出。一般地，设 \(q_{1}, \ldots, q_{l}\) 为包含于 p 中的 B 的极小素理想。使用 Nullstellensatz 证明，存在 \(V(p)\) 的稠密开集 U，使得对于属于 U 的每个极大理想 m，包含于 m 中的极小素理想都包含于 U。使用 c) 证明，缩小 U 后还可以假设对于每个 \(p' \in U\)，有 \(\dim(\mathrm{P}_{\mathrm{B/A}}^{\infty}[\mathfrak{p}]) = \dim(\mathrm{P}_{\mathrm{B/A}}^{\infty}[\mathfrak{p}'])\)。得出结论。）

\begin{enumerate}
\def\labelenumi{\alph{enumi})}
\setcounter{enumi}{5}
\tightlist
\item
  假设 B 为有限生成的 A-代数。证明函数 \(\mathfrak{p} \mapsto \mathrm{d}(\mathfrak{p}) = \dim_{\mathfrak{p}} (\mathrm{B} \otimes_{\mathbf{A}} \kappa(\rho^{-1}(\mathfrak{p})))\) 是上半连续的。（当 B 有限表示时，使用 \(e\) 和 \(c\)。一般地，B 具有形式 \(\mathrm{A}[\mathrm{T}_1, ..., \mathrm{T}_n]/\mathfrak{J}\)。令 \(\mathfrak{J}_{\alpha}\) 为 \(\mathrm{A}[\mathrm{T}_1, ..., \mathrm{T}_n]\) 中包含于 \(\mathfrak{J}\) 的有限生成理想族，并按包含关系排序，置 \(\mathrm{B}_{\alpha} = \mathrm{A}[\mathrm{T}_1, ..., \mathrm{T}_n]/\mathfrak{J}_{\alpha}\)。对于每个 \(\alpha\)，有 \(\operatorname{Spec}(\mathbf{B}) \subset \operatorname{Spec}(\mathrm{B}_{\alpha}) \subset \operatorname{Spec}(\mathrm{A}[\mathrm{T}_1, ..., \mathrm{T}_n])\)，且 \(\operatorname{Spec}(\mathbf{B}) = \bigcap \operatorname{Spec}(\mathrm{B}_{\alpha})\)。对于每个 \(\mathfrak{p} \in \operatorname{Spec}(\mathbf{B})\)，置
\end{enumerate}

\[
\mathrm{d} _ {\alpha} (\mathfrak {p}) = \dim_ {\mathfrak {p}} \left(\mathrm{B} _ {\alpha} \otimes \kappa \left(\rho^ {- 1} (\mathfrak {p})\right)\right).
\]

  证明有 \(d_{\alpha}(p) \geqslant d(p)\) 及 \(d(p) = \inf_{\alpha} d_{\alpha}(p)\)，这是根据 \(\text{Spec}(\kappa(\rho^{-1}(p))[T_1, \ldots, T_n])\) 的诺特性。推出 \(p \mapsto d(p)\) 是一个递减滤族上半连续函数的下确界。）

\begin{enumerate}
\def\labelenumi{\arabic{enumi})}
\setcounter{enumi}{10}
\tightlist
\item
  设 \(\rho: \mathbf{A} \to \mathbf{B}\) 为环同态，\(^a\rho: \operatorname{Spec}(\mathbf{B}) \to \operatorname{Spec}(\mathbf{A})\) 为相应映射。
\end{enumerate}

\begin{enumerate}
\def\labelenumi{\alph{enumi})}
\item
  证明对于每个 \(\mathfrak{q} \in \operatorname{Spec}(\mathbf{A})\)，\(\operatorname{Spec}(\mathbf{B} \otimes_{\mathbf{A}} \kappa(\mathfrak{q}))\) 可与 \(^a\rho^{-1}(\mathfrak{q})\) 识别。
\item
  假设 B 为有限生成的 A-代数。证明 p 是其纤维 \(^a\rho^{-1}(^a\rho(p))\) 中的孤立点，当且仅当 \(\dim_p(B \otimes_A \kappa(^a\rho(p)))\) 为零。
\item
  从习题 10 推出，在 \(b\) 的假设下，\(\operatorname{Spec}(\mathbf{B})\) 中在各自纤维内孤立的点集是 \(\operatorname{Spec}(\mathbf{B})\) 的开子集。
\end{enumerate}

\begin{enumerate}
\def\labelenumi{\arabic{enumi})}
\setcounter{enumi}{11}
\tightlist
\item
  设 H 为有限生成的分次代数，\(H_{0}\) 是一个域，\((\xi_{a})_{a\in\mathbf{A}}\) 是生成 H 的齐次元素。假设存在严格正整数族 \((d_{i})_{i\in\mathbf{I}}\)，使得 \(P_{H}=\prod(1-T^{d_{i}})^{-1}\)。以 \(\delta_{a}>0\) 表示 \(\xi_{a}\) 的次数。
\end{enumerate}

\begin{enumerate}
\def\labelenumi{\alph{enumi})}
\item
  证明存在单射 \(\varphi: I \to A\)，使得对于每个 \(i\)，\(d_i\) 整除 \(\delta_{\varphi(i)}\)。（可以使用 § 4, no. 2, 定理 1。）
\item
  证明有 \(\inf_{a \in A} \delta_a \leqslant \inf_{i \in I} d_i\)。
\item
  由此推出，若各 \(\delta_{a}\) 彼此相等，则分次代数 H 是正则的。
\item
  假设 A 有两个元素 \(a\) 和 \(a'\)，使得 \(\delta_{a} < \delta_{a'}\)，且 \(\delta_{a'}\) 不是 \(\delta_{a}\) 的倍数。证明分次代数 H 是正则的。
\end{enumerate}

